\UseRawInputEncoding
\documentclass[conference]{IEEEtran}
\IEEEoverridecommandlockouts
\usepackage{cite}
\usepackage{amsmath,amssymb,amsfonts}
\usepackage{algorithmic}
\usepackage{graphicx}
\usepackage{afterpage}
\usepackage{textcomp}
\usepackage{xcolor}
\usepackage{url}
\usepackage{listings}
\usepackage[utf8]{inputenc}
\usepackage{minted}
\usepackage{float}
\usepackage{subcaption}

\lstset{
  language=Python,
  basicstyle=\ttfamily\footnotesize,
  breaklines=true,
  breakatwhitespace=true,
  columns=fullflexible,
  frame=single
}

\setminted{
  fontsize=\footnotesize,
  breaklines=true,
  frame=single,
  linenos=true,
  tabsize=2,
  autogobble=true
}

\def\BibTeX{{\rm B\kern-.05em{\sc i\kern-.025em b}\kern-.08em
    T\kern-.1667em\lower.7ex\hbox{E}\kern-.125emX}}

\begin{document}

\title{Servidor HTTP}

\author{\IEEEauthorblockN{David Acuña López} \IEEEauthorblockA{\textit{Escuela de Ingeniería en Computación} \\ \textit{Tecnológico de Costa Rica}\\ Cartago, Costa Rica \\ \texttt{rodolfoide69@estudiantec.cr}} \and \IEEEauthorblockN{Raúl Sanabria Marroquín} \IEEEauthorblockA{\textit{Escuela de Ingeniería en Computación} \\ \textit{Tecnológico de Costa Rica}\\ Cartago, Costa Rica \\ \texttt{raulsanabria@estudiantec.cr}} }

\maketitle

\begin{abstract}
Se implementó un servidor HTTP en Rust con routing, colas y pools de trabajo diferenciados (básico, CPU e IO). El sistema expone \textit{handlers} de carga sintética (CPU-bound e IO-bound), utilitarios y de archivos, junto con endpoints de observabilidad (\texttt{/status}, \texttt{/metrics}). Se compara el impacto de tamaño de pool, profundidad de cola y política de planificación sobre latencias (p50/p95/p99) y throughput. Los resultados muestran las compensaciones entre paralelismo y contención, y cuantifican el efecto de \textit{backpressure} en escenarios con \textit{workloads} heterogéneos.
\end{abstract}

\begin{IEEEkeywords}
Rust, HTTP, colas, pools de hilos, CPU-bound, IO-bound, latencia p50/p95/p99, escalabilidad, planificación.
\end{IEEEkeywords}

\section{Introducción}
\begin{itemize}
  \item {Este trabajo presenta un servidor HTTP en Rust} que integra:
    \begin{itemize}
      \item {Ruteo liviano por tabla.}
      \item {Tres pools con colas separadas:} \emph{básico}, \emph{CPU} e \emph{IO}.
      \item {Handlers representativos:}
        \begin{itemize}
          \item \emph{CPU–intensivos:} \texttt{/isprime}, \texttt{/factor}, \texttt{/pi}, \texttt{/mandelbrot}, \texttt{/matrixmul}.
          \item \emph{IO–intensivos:} \texttt{/sortfile}, \texttt{/wordcount}, \texttt{/grep}, \texttt{/compress}, \texttt{/hashfile}.
        \end{itemize}
      \item {Utilitarios:} \texttt{/reverse}, \texttt{/toupper}, \texttt{/random}, \texttt{/hash}, \texttt{/sleep}, \texttt{/createfile}, \texttt{/deletefile}.
      \item {Observabilidad:} \texttt{/status}, \texttt{/metrics}.
    \end{itemize}
\end{itemize}

\noindent\textbf{Objetivo.} Estudiar latencias (\emph{p50/p95/p99}) y escalabilidad al variar el tamaño de los pools, la profundidad de las colas y la carga.

\section{Marco Teórico}
\subsection{CPU-bound vs IO-bound}
\textbf{CPU-bound} consume ciclos de CPU: factorización, primalidad, $\pi$ (spigot/Machin), Mandelbrot, multiplicación de matrices. \textbf{IO-bound} depende de disco/FS: lectura/escritura, compresión, hashing de archivos.

\subsection{Pools y colas}
Separar colas y pools por tipo de carga evita head-of-line blocking. Profundidad de cola limita \emph{backlog}; políticas de despacho (FCFS por cola) y \emph{timeouts} controlan fairness y acotación de espera.

\subsection{Métricas de latencia y throughput}
Se reportan p50/p95/p99, tiempo de servicio, y throughput en req/s. El cálculo de percentiles sigue la práctica de ordenar muestras por endpoint y extraer cuantiles empíricos.

\section{Diseño e Implementación}
\subsection{Arquitectura}
\textbf{Core:} \texttt{AppState} mantiene config, router, métricas, pools y almacén de jobs.  
\textbf{Pools:} \texttt{basic}, \texttt{cpu}, \texttt{io} con colas y workers configurables.  
\textbf{Observabilidad:}
\texttt{/status} expone configuración y snapshot de colas; \texttt{/metrics} entrega contadores, latencias (avg/p50/p95/p99) y throughput.

\subsection{Ruteo y Handlers}
Cada handler sigue \texttt{(state:\&Shared, req:\&Request) -> (status, ctype, body)}. A modo de ejemplo:
\begin{itemize}
  \item \textbf{CPU-bound:} \texttt{/isprime}, \texttt{/factor}, \texttt{/pi?digits=D}, \texttt{/mandelbrot}, \texttt{/matrixmul?size=N\&seed=S}.
  \item \textbf{IO-bound:} \texttt{/sortfile}, \texttt{/wordcount}, \texttt{/grep}, \texttt{/compress?codec=gzip|xz}, \texttt{/hashfile?algo=sha256}.
  \item \textbf{Utilitarios:} \texttt{/reverse}, \texttt{/toupper}, \texttt{/random}, \texttt{/hash}, \texttt{/sleep}, \texttt{/createfile}, \texttt{/deletefile}.
\end{itemize}

\subsection{Configuración}
\textbf{Parámetros:} \texttt{workers\_basic, workers\_cpu, workers\_io}, \texttt{queue\_basic, queue\_cpu, queue\_io}, \texttt{timeout\_cpu\_ms, timeout\_io\_ms}.  
Estos se reflejan en \texttt{/status} y afectan el \emph{scheduling} efectivo.

\section{Estrategia de Pruebas}
\subsection{Metodología de carga}
Se generan series de requests por endpoint con lotes de tamaño $T$ y pausa controlada. Para IO-bound se crean archivos en \texttt{data/}. Para CPU-bound se varía `digits`, `size`, `n`, etc.

\subsection{Matriz experimental}
\begin{table}[H]
\centering
\caption{Diseño factorial resumido (TODO: ajustar a corridas reales)}
\begin{tabular}{lccc}
\hline
Factor & Niveles & Valores ejemplo & Afecta \\
\hline
Workers CPU & 3 & \{1,2,4\} & p95/p99 CPU \\
Workers IO  & 3 & \{1,2,4\} & p95/p99 IO \\
Queue depth & 2 & \{64,256\} & cola/backlog \\
Carga        & 3 & \{10,50,100 rps\} & colas \\
\hline
\end{tabular}
\end{table}

\subsection{Recolección de métricas}
Cada $x$ segundos se consulta \texttt{/metrics} y se almacena JSON con:
\texttt{accepted, handled, errors, latency\_ms\{avg,p50,p95,p99\}, throughput, queues}.

\section{Resultados}
% TODO: sustituir por figuras/tablas reales.
\subsection{Latencias por endpoint}
\begin{figure}[H]
\centering
\includegraphics[width=0.9\linewidth]{fig_p95_por_endpoint.png}
\caption{p95 por endpoint bajo configuración base (TODO).}
\end{figure}

\subsection{Efecto del tamaño de pool}
\begin{table}[H]
\centering
\caption{p50/p95/p99 vs workers\_cpu (carga fija) (TODO)}
\begin{tabular}{lcccc}
\hline
workers\_cpu & p50 & p95 & p99 & rps \\
\hline
1 & -- & -- & -- & -- \\
2 & -- & -- & -- & -- \\
4 & -- & -- & -- & -- \\
\hline
\end{tabular}
\end{table}

\subsection{Profundidad de cola y backpressure}
\begin{figure}[H]
\centering
\includegraphics[width=0.9\linewidth]{fig_backpressure.png}
\caption{Crecimiento de pendientes en cola vs latencia (TODO).}
\end{figure}

\section{Discusión}
(1) \textbf{CPU-bound} escala con \#workers hasta saturar cores; p99 mejora al aumentar workers\_cpu pero se degrada si hay \textit{oversubscription}.  
(2) \textbf{IO-bound} se beneficia de mayor concurrencia, pero colas profundas elevan p95/p99 por espera en cola (HoL blocking mitigado al separar colas).  
(3) \textbf{Queue depth} alto amortigua ráfagas pero incrementa latencias bajo alta carga sostenida.  
(4) \textbf{Política de planificación:} FCFS por cola es suficiente en esta arquitectura por separación de clases de trabajo; preempción a nivel de handler no aporta en IO-bound, sí en CPU-bound largo si se quisiera fairness intra-cola.

\section{Conclusiones}
Separar pools/colas por tipo de carga controla head-of-line blocking y mejora p95/p99. Aumentar workers mejora throughput hasta el punto de saturación del recurso (CPU o disco) y luego aumenta colisiones y p99. La profundidad de cola es un \emph{trade-off} entre absorción de ráfagas y latencia tail. La instrumentación con \texttt{/metrics} permite ajustar configuración en función de objetivos (SLOs) por endpoint.

\section{Referencias}
\begin{thebibliography}{9}
\bibitem{rust_book} The Rust Programming Language. \url{https://doc.rust-lang.org/}
\bibitem{tokio} Tokio runtime (conceptos de IO asíncrono). \url{https://tokio.rs/}
\bibitem{tail_latency} Dean, Barroso. The Tail at Scale. \url{https://research.google/pubs/the-tail-at-scale/}
\end{thebibliography}

\end{document}
